Class imbalance is overwhelmingly handled by oversampling, yet SMOTE-family methods
typically trade majority-class accuracy for minority-class recall, and it is rarely
clear \emph{why} oversampling helps. We introduce a \emph{remedy decomposition} of
the misclassified minority: an out-of-bag error triage (ensemble confidence and
local class representation, validated against an aleatoric--epistemic uncertainty
decomposition) is combined with a threshold-recoverability test, assigning each
minority error to the remedy that would fix it: \emph{threshold-recoverable} (its
score already clears the tuned decision threshold; only the default cutoff is
wrong), \emph{data-reducible} (more same-class data would help), or
\emph{irreducible} (genuine class overlap). Across 79 real datasets the deficit is
predominantly threshold-recoverable (mean $68\%$; $80\%$ at imbalance ratio $>3$);
only $\sim$6\% is data-reducible and $\sim$17\% irreducible. This explains a known
but unexplained regularity: SMOTE-family resampling does not improve ranking (AUC
is unchanged) and a threshold move reproduces its benefit without synthetic
generation; single undertrained neural networks, whose ranking imbalance itself
corrupts, are the exception. On controlled mixtures with a
\emph{known} Bayes boundary, standard SMOTE places an overlap-dependent fraction of
synthetic points \emph{across} that boundary, and a short proposition shows why
interpolating boundary seeds crosses into majority-posterior space. Once every strategy receives
the same out-of-fold threshold tuning, the default-threshold gap between
non-generative and generative families collapses to within $\pm1$\,pp across three
base learners (against $+6.4$\,pp from tuning the threshold alone): the SMOTE
family's apparent edge is an operating-point effect, not a discrimination gain. The
practical conclusion: diagnose the deficit; for most of it, move the decision
rather than oversample.
